\section*{Bab \ref{ch:b2:angiosperm-reproduction} --- Perkembangbiakan Seksual Tumbuhan Berbunga}

\begin{solution}{exo:b2:angiosperm-reproduction:1}
Daun kelopak (kelopaknya), daun mahkota (mahkotanya), \omterm{def:b2:angiosperm-reproduction:flower}{benang sari}
(androsium), \omterm{def:b2:angiosperm-reproduction:flower}{daun buah} (ginesium). \omterm{def:b2:angiosperm-reproduction:flower}{Benang sari}: tangkai sari dan kepala
sari. \omterm{def:b2:angiosperm-reproduction:flower}{Daun buah}: \omterm{def:b2:angiosperm-reproduction:flower}{kepala putik}, tangkai putik, \omterm{def:b2:angiosperm-reproduction:flower}{bakal buah} berisi \omterm{def:b2:plant-life-cycles:heterospory}{bakal
biji}. Semuanya \omterm{def:b2:plant-life-cycles:cycles}{sporofit} yang \omterm{def:b2:meiosis-heredity:meiosis}{diploid}; hanya \omterm{prop:b2:angiosperm-reproduction:pollen}{butir serbuk sari} (di dalam
kepala sarinya) dan \omterm{prop:b2:angiosperm-reproduction:embryosac}{kantung lembaga} (di dalam \omterm{def:b2:plant-life-cycles:heterospory}{bakal bijinya}) yang
merupakan \omterm{def:b2:plant-life-cycles:cycles}{gametofit} \omterm{def:b2:meiosis-heredity:meiosis}{haploid}.
\end{solution}

\begin{solution}{exo:b2:angiosperm-reproduction:2}
\omterm{def:b2:plant-life-cycles:heterospory}{Mikrospora} $n$; sel vegetatif $n$; sperma $n$; sel telur $n$; \omterm{prop:b2:angiosperm-reproduction:embryosac}{inti
kutub} $n$; zigot $2n$; endosperma $3n$; \omterm{def:b2:angiosperm-reproduction:seed}{kulit biji} $2n$ (integumen dari
induknya); dinding \omterm{def:b2:angiosperm-reproduction:seed}{buah} $2n$ (\omterm{def:b2:angiosperm-reproduction:flower}{bakal buah} dari induknya).
\end{solution}

\begin{solution}{exo:b2:angiosperm-reproduction:3}
Angin (rumput, ek, birka, hazel, sendokan): \omterm{def:b2:angiosperm-reproduction:flower}{bunga} hijau kecil, tanpa
daun mahkota, wangi, maupun \omterm{def:b2:angiosperm-reproduction:pollination}{nektar}, kepala sari yang bergantung pada
tangkai sari yang panjang, \omterm{def:b2:angiosperm-reproduction:flower}{kepala putik} berbulu, \omterm{def:b2:plant-life-cycles:heterospory}{serbuk sari} kering dan
licin dalam jumlah luar biasa banyak, dan berbunga sebelum daunnya
muncul. Lebah (semanggi, salvia, digitalis, lavender, apel): daun
mahkota berwarna yang kerap berpenuntun ultraungu, landasan pendaratan,
wangi, \omterm{def:b2:angiosperm-reproduction:pollination}{nektar}, \omterm{def:b2:plant-life-cycles:heterospory}{serbuk sari} lengket berpahat dalam jumlah sedang, dan
berbunga di siang hari.
\end{solution}

\begin{solution}{exo:b2:angiosperm-reproduction:4}
Kedua sel sperma satu \omterm{thm:b2:angiosperm-reproduction:tube}{buluh serbuk sari} melebur, satu dengan sel
telurnya (zigot, $2n$, lembaganya) dan satu dengan sel pusatnya (inti
endosperma, $3n$, simpanan makanannya). Endospermanya baru mulai tumbuh
sesudah pembuahan, sehingga tumbuhannya hanya membekali \omterm{def:b2:plant-life-cycles:heterospory}{bakal biji} yang
memuat sebuah lembaga, tidak seperti gimnosperma, yang membangun
simpanan makanannya lebih dahulu.
\end{solution}

\begin{solution}{exo:b2:angiosperm-reproduction:5}
$25/1.2 = \qty{21}{h}$. Rambut yang muncul pada hari ke-5, ke-6, dan
ke-7 berjumpa \omterm{def:b2:plant-life-cycles:heterospory}{serbuk sari}; rambut hari ke-8 sampai ke-12 tidak: $3/8$
rambutnya terserbuki, dan lima perdelapan tongkolnya hampa.
\end{solution}

\begin{solution}{exo:b2:angiosperm-reproduction:6}
$S_1S_2\times S_1S_2$: 0. $S_1S_2\times S_1S_3$: $\tfrac12$ (hanya
\omterm{def:b2:plant-life-cycles:heterospory}{serbuk sari} $S_3$ yang tumbuh). $S_1S_2\times S_3S_4$: semuanya.
Keturunan persilangan yang kedua: sel telur $S_1$ atau $S_2$, sperma
$S_3$: $S_1S_3$ dan $S_2S_3$, masing-masing separuh.
\end{solution}

\begin{solution}{exo:b2:angiosperm-reproduction:7}
Sebuah tumbuhan mengusung 2 dari $n$ \omterm{def:b2:meiosis-heredity:vocabulary}{alelnya} dan menolak \omterm{def:b2:plant-life-cycles:heterospory}{serbuk sari}
yang mengusung salah satunya: bagian yang serasi $(n - 2)/n$. Sebuah
\omterm{def:b2:meiosis-heredity:vocabulary}{alel} baru tidak ditolak putik mana pun, sehingga \omterm{def:b2:plant-life-cycles:heterospory}{serbuk sarinya}
mengayahi lebih banyak \omterm{def:b2:angiosperm-reproduction:seed}{biji} daripada rata-rata dan ia menyebar sampai
sama lazimnya dengan yang lain; keunggulan yang sama melindungi tiap
\omterm{def:b2:meiosis-heredity:vocabulary}{alel} yang langka dari kepunahan. Maka seleksi menimbun \omterm{def:b2:meiosis-heredity:vocabulary}{alel}, dan
populasi alami spesies yang takserasi sendiri mengusung puluhan \omterm{def:b2:meiosis-heredity:vocabulary}{alel}.
\end{solution}

\begin{solution}{exo:b2:angiosperm-reproduction:8}
Kedua \omterm{prop:b2:angiosperm-reproduction:embryosac}{inti kutubnya} adalah salinan satu \omterm{def:b2:plant-life-cycles:heterospory}{megaspora}, sehingga sel
pusatnya $YY$ atau $yy$ (masing-masing separuh); spermanya $Y$ atau $y$
(masing-masing separuh): endospermanya $YYY$, $YYy$, $Yyy$, $yyy$
masing-masing $\tfrac14$; tiga perempatnya kuning. Kalau diserbuki
$yy$: $YYy$ dan $yyy$, masing-masing separuh --- separuh bulirnya
kuning.
\end{solution}

\begin{solution}{exo:b2:angiosperm-reproduction:9}
Bulir jelai dipotong menjadi dua; belahan yang berlembaga mencerna
patinya, belahan yang tak berlembaga tidak, kecuali kalau giberelin
ditambahkan, dan barulah ia mencernanya. (a) Lembaganya adalah sumber
isyaratnya, bukan sumber enzimnya. (b) Giberelinnya adalah isyaratnya,
dan sudah cukup dengan sendirinya. (c) Amilasenya dibuat lapisan
aleuron endospermanya, yang menanggapi hormonnya. Belahan yang tak
berlembaga memisahkan isyaratnya dari tanggapannya: tanpanya orang
tidak dapat mengetahui apakah lembaganya sendiri yang mencerna patinya.
\end{solution}

\begin{solution}{exo:b2:angiosperm-reproduction:10}
Angin: $10^{6}\times\qty{1}{\micro g} = \qty{1}{g}$ tiap \omterm{def:b2:plant-life-cycles:heterospory}{bakal biji}.
Hewan: \qty{1}{mg} \omterm{def:b2:plant-life-cycles:heterospory}{serbuk sari} ditambah \qty{0.5}{mg} gula
$= \qty{1.5}{mg}$ tiap \omterm{def:b2:plant-life-cycles:heterospory}{bakal biji} --- kira-kira tujuh ratus kali lebih
murah. \omterm{def:b2:angiosperm-reproduction:pollination}{Penyerbukan} angin tidak membutuhkan pasangan: ia berhasil pada
awal musim semi sebelum serangga terbang, di tempat yang dingin,
berangin, dan terbuka, pada malam hari dan di dalam hujan, di dalam
rumpun rapat satu spesies, dan ia tidak dapat dicurangi pencuri \omterm{def:b2:angiosperm-reproduction:pollination}{nektar}
maupun ditinggalkan penyerbuk yang punah.
\end{solution}

\begin{solution}{exo:b2:angiosperm-reproduction:11}
\omterm{def:b2:angiosperm-reproduction:seed}{Biji} yang jatuh di kaki tetuanya mendarat di dalam naungannya, di
antara akarnya, dan di tempat pemangsa \omterm{def:b2:angiosperm-reproduction:seed}{biji} serta patogen khasnya
berkumpul; hampir semuanya mati. Seekor burung mengangkut sebutir \omterm{def:b2:angiosperm-reproduction:seed}{biji}
berkilo-kilometer (\qty{12}{km} dalam dua puluh menit) lalu
menjatuhkannya, bersama pupuk, di sebuah pagar tanaman atau sebuah
tanah terbuka; sekalipun hanya satu \omterm{def:b2:angiosperm-reproduction:seed}{biji} dari seratus yang mendarat di
tempat yang baik, itu lebih banyak daripada seluruh \omterm{def:b2:angiosperm-reproduction:seed}{biji} di bawah
pohonnya, dan populasinya menyebar pada kelajuan burung, bukan pada
kelajuan \omterm{def:b2:angiosperm-reproduction:seed}{buah} yang jatuh. Dagingnya adalah harga tiketnya.
\end{solution}

\begin{solution}{exo:b2:angiosperm-reproduction:12}
Pengurungannya membuat \omterm{def:b2:plant-life-cycles:heterospory}{serbuk sarinya} harus tumbuh menembus jaringan
induknya, dan itu membiarkan putiknya mengujinya (\omterm{prop:b2:angiosperm-reproduction:selfing}{ketakserasian
sendiri}, penolakan spesies asing), membiarkan banyak buluh bersaing
sehingga yang tercepat mengayahi \omterm{def:b2:angiosperm-reproduction:seed}{bijinya}, melindungi \omterm{def:b2:plant-life-cycles:heterospory}{bakal bijinya} dari
kekeringan dan dari pemakan tumbuhan, serta mengubah \omterm{def:b2:angiosperm-reproduction:flower}{bakal buahnya}
menjadi \omterm{def:b2:angiosperm-reproduction:seed}{buah} yang merekrut hewan untuk penyebarannya; \omterm{prop:b2:angiosperm-reproduction:double}{pembuahan
gandanya} lalu hanya membekali \omterm{def:b2:plant-life-cycles:heterospory}{bakal biji} yang telah dibuahi. Ongkosnya:
sebutir \omterm{def:b2:plant-life-cycles:heterospory}{serbuk sari} harus mencapai sebuah \omterm{def:b2:angiosperm-reproduction:flower}{kepala putik} alih-alih \omterm{def:b2:plant-life-cycles:heterospory}{bakal
bijinya} sendiri, sehingga tangkai putik, sebuah buluh, dan
penuntunannya dibutuhkan, dan \omterm{def:b2:angiosperm-reproduction:seed}{buahnya} adalah penanaman yang berat.
Neracanya memihak angiosperma begitu kuat sehingga mereka beranjak dari
tiada menjadi sembilan persepuluh spesies tumbuhan dalam seratus juta
tahun.
\end{solution}

\begin{solution}{pb:b2:angiosperm-reproduction:1}
\textbf{1.} Pada A ($S_1S_2$): \omterm{def:b2:plant-life-cycles:heterospory}{serbuk sari} A $f = 0$; B ($S_1$, $S_3$) $f =
\tfrac12$; C ($S_3$, $S_4$) $f = 1$.
\textbf{2.} Pada B ($S_1S_3$): A $\tfrac12$, B 0, C $\tfrac12$. Pada C
($S_3S_4$): A 1, B $\tfrac12$, C 0.
\textbf{3.} 50 butir yang serasi; $0.1\times 50 = 5$ \omterm{def:b2:plant-life-cycles:heterospory}{bakal biji}.
\textbf{4.} Dari C: 100 yang serasi, $0.1\times 100 = 10$, yakni
kesepuluh \omterm{def:b2:plant-life-cycles:heterospory}{bakal bijinya}; dari A: tidak ada.
\textbf{5.} \omterm{def:b2:angiosperm-reproduction:seed}{Buah} dari kunjungan C (10 \omterm{def:b2:angiosperm-reproduction:seed}{biji}) dipertahankan; dari B (5
\omterm{def:b2:angiosperm-reproduction:seed}{biji}) pas dipertahankan; dari A dijatuhkan. Sebuah kebun berisi A saja
tidak berbuah sama sekali.
\textbf{6.} Separuh \omterm{def:b2:angiosperm-reproduction:flower}{bunganya} menjadikan \omterm{def:b2:angiosperm-reproduction:seed}{buah} (yang dikunjungi dari C):
2500 \omterm{def:b2:angiosperm-reproduction:seed}{buah} tiap pohon, sepuluh kali 250 yang dibutuhkan --- panen penuh,
dan pohonnya akan menjatuhkan kelebihannya.
\textbf{7.} Ia melepaskan \omterm{def:b2:plant-life-cycles:heterospory}{serbuk sari} yang serasi sepanjang seluruh
musim berbunga tiap varietasnya; tetapi kalau ia berbagi kedua \omterm{def:b2:meiosis-heredity:vocabulary}{alel} $S$
dengan sebuah varietas, ia tidak menyerbuki apa pun padanya.
\textbf{8.} $8\times 10^{7}/(7\times 86400) = 132$ butir tiap meter
persegi tiap sekon.
\textbf{9.} $132/0.25 = 530$ butir tiap meter kubik.
\textbf{10.} $132\times 10^{-4} = 0.013$ tiap sekon: satu butir tiap
\qty{76}{s}.
\textbf{11.} $0.013\times 86400 \approx 1100$ butir sehari; buluh
pertama yang mencapai \omterm{def:b2:plant-life-cycles:heterospory}{bakal bijinya} membuahinya, dan \omterm{def:b2:plant-life-cycles:heterospory}{bakal bijinya} lalu
tertutup bagi yang lain.
\textbf{12.} \qty{20}{h} sesudah \omterm{def:b2:angiosperm-reproduction:pollination}{penyerbukan}.
\textbf{13.} \omterm{def:b2:plant-life-cycles:heterospory}{Serbuk sarinya} pada hari ke-0 sampai ke-7, rambutnya sejak
hari ke-5: hanya hari ke-5, ke-6, dan ke-7 yang bertindih --- $3/8$
rambutnya, dan rambut yang muncul sesudah hari ke-7 tidak memperoleh
apa pun; tongkolnya sebagian besar hampa. Kekeringan pada saat
berambut adalah penyebab klasik gagalnya panen jagung.
\textbf{14.} Lembaganya $Yy$; endospermanya $YYy$.
\textbf{15.} Kuning, karena $Y$ \omterm{def:b2:meiosis-heredity:vocabulary}{dominan}: ciri tetua \omterm{def:b2:plant-life-cycles:heterospory}{serbuk sarinya}
tampak pada \omterm{def:b2:angiosperm-reproduction:seed}{biji} di atas tumbuhan induknya --- yakni xenia.
\textbf{16.} Sel pusatnya $YY$ atau $yy$: endospermanya $YYy$
(separuhnya) dan $yyy$ (separuhnya).
\textbf{17.} $YYy$: 20 satuan; $Yyy$: 10 satuan; sebuah $yyY$ dengan
$Y$ dari ayahnya adalah \omterm{def:b2:meiosis-heredity:vocabulary}{genotipe} yang sama $Yyy$, yaitu 10 satuan. Dua
takaran dari ayahnya tidak dapat ada: satu sperma membuahi sel
pusatnya, dan endospermanya selalu dua genom dari induknya berbanding
satu dari ayahnya.
\textbf{18.} Ia adalah simpanan patinya; $0.3/0.33 = \qty{91}{\%}$
bulirnya.
\textbf{19.} Simpanannya baru dibangun sesudah pembuahan, sehingga tidak
ada sumber daya yang masuk ke \omterm{def:b2:plant-life-cycles:heterospory}{bakal biji} yang tidak mengusung lembaga;
\omterm{def:b2:plant-life-cycles:cycles}{gametofit} gimnosperma dibangun lebih dahulu dan terbuang kalau
\omterm{def:b2:angiosperm-reproduction:pollination}{penyerbukannya} gagal.
\textbf{20.} $0.9\times 600 = 540$ bulir tiap tongkol dan tiap
tumbuhan; $4320$ tiap meter persegi; $4320\times 0.3 = \qty{1.3}{kg}$
endosperma tiap meter persegi.
\textbf{21.} \qty{13}{t/ha}.
\textbf{22.} Karbon bulirnya $0.5\times 1.2 = \qty{0.6}{kg/m^2}$;
karbon endospermanya $0.45\times 1.3 = \qty{0.58}{kg/m^2}$: sepadan
(sisanya yang kecil adalah lembaga dan kulitnya).
\textbf{23.} $8\times 10^{7}/4320 \approx \num{18500}$ butir tiap
bulir.
\textbf{24.} Awan \omterm{def:b2:plant-life-cycles:heterospory}{serbuk sarinya} sendiri tipis dan tertiup pergi, dan
malainya melepaskan \omterm{def:b2:plant-life-cycles:heterospory}{serbuk sari} sebagian besar sebelum rambutnya
sendiri muncul, sehingga banyak rambutnya tidak pernah terserbuki dan
tongkolnya berlubang-lubang.
\textbf{25.} Bidang A: separuh \omterm{def:b2:angiosperm-reproduction:flower}{bunganya} menjadikan \omterm{def:b2:angiosperm-reproduction:seed}{buah}, 2500 tiap
pohon, yaitu panen penuh. Ladangnya: 4320 bulir tiap meter persegi,
yaitu \qty{13}{t/ha}.
\end{solution}
